\chapter{Results and Evaluation}
\label{ch:results}

This chapter reports the measured outcome of the design and verification work. It first
fixes the exact baseline and counting method so that every number below is reproducible
(\cref{sec:results-baseline}), then presents the implementation metrics
(\cref{sec:impl-metrics}), the regression outcome (\cref{sec:regression}), waveform
evidence (\cref{sec:waveforms}), and the functional-coverage result
(\cref{sec:funccov}). It closes by comparing the controller with the CHIPS Alliance
reference (\cref{sec:ref-comparison}) and by evaluating the design against the requirements
defined in \cref{ch:background-requirements} (\cref{sec:req-eval}). Interpretation of the remaining
limitations and the roadmap is deferred to \cref{ch:conclusion}.

\section{Verified baseline and measurement method}
\label{sec:results-baseline}

All figures in this chapter come from a single verified snapshot of the repository, run on
one simulator with one seed. \Cref{tab:results-baseline} records that baseline. Fixing it
matters because the reported counts are volatile: they change as sequences and RTL evolve,
so a result is meaningful only against the commit and tool that produced it.

\begin{table}[!htb]\centering
  \caption{Verified baseline for all results in this chapter.}
  \label{tab:results-baseline}
  \small
  \begin{tabular}{@{}>{\raggedright\arraybackslash}p{0.34\textwidth}
                    >{\raggedright\arraybackslash}p{0.58\textwidth}@{}}
    \toprule
    \textbf{Item} & \textbf{Value} \\
    \midrule
    Repository branch & \texttt{feature/functional-coverage-model} \\
    Commit & \texttt{047dfdc} (``Add functional coverage model'') \\
    Simulator & Cadence Xcelium \texttt{xrun} 18.03-s001 \\
    Verification library & UVM~1.2 (\texttt{CDNS-1.2}) \\
    Seed policy & fixed \texttt{-svseed 1}, single pass per sequence \\
    Coverage flow & \texttt{-coverage all} $\rightarrow$ per-test UCIS databases, merged and
      summarised by the project \texttt{ucd2json}/\texttt{extract\_coverage.py} tools \\
    Reference design & CHIPS Alliance \texttt{i3c-core} @ \texttt{a9a8ef9b} \\
    \bottomrule
  \end{tabular}
\end{table}

\paragraph{Counting rules.} Two selection rules are applied consistently. \emph{Line counts}
are taken with \texttt{wc -l} over \texttt{*.sv} files: the project design is the
\texttt{src/rtl} tree; the reference design is its \texttt{src} tree. Where an RTL-only
figure is quoted, the reference's PeakRDL-generated UVM/coverage helper files
(\texttt{*\_uvm.sv}, \texttt{*\_sample.svh}, \texttt{*\_covergroups.svh}) are excluded, since
they are verification artefacts rather than synthesisable logic. \emph{Runnable sequences}
are the virtual sequences that drive a named scenario; the four abstract
\texttt{*\_base\_vseq} classes and the non-power-of-two FIFO elaboration guard are excluded,
leaving \RunnableVseqCount{} runnable scenario sequences. These same \RunnableVseqCount{}
sequences form both the regression population and the coverage population.

\section{Quantitative implementation metrics}
\label{sec:impl-metrics}

The controller RTL is 5121 lines across 18 SystemVerilog files.
\Cref{tab:impl-subsystem} breaks this down by subsystem. The command-flow engine dominates:
\texttt{flow\_active.sv} alone is 1794 lines, consistent with its role as the central
\FlowStateCount{}-state protocol FSM. Everything else --- CSR/DAT, HCI queues, PHY, and the
top level --- is comparatively small, which is the intended shape of a controller whose
complexity is concentrated in one auditable state machine.

\begin{table}[!htb]\centering
  \caption{Per-subsystem line count of the controller RTL (\texttt{src/rtl}).}
  \label{tab:impl-subsystem}
  \small
  \begin{tabular}{@{}>{\raggedright\arraybackslash}p{0.46\textwidth} r r@{}}
    \toprule
    \textbf{Subsystem} & \textbf{Files} & \textbf{Lines} \\
    \midrule
    Control FSMs and datapath (\texttt{ctrl}) & 12 & 3997 \\
    \quad of which \texttt{flow\_active.sv} & 1 & 1794 \\
    CSR register file + 32-entry DAT (\texttt{csr}) & 1 & 368 \\
    HCI queues (\texttt{hci}) & 2 & 241 \\
    PHY (\texttt{phy}) & 1 & 55 \\
    Top level + shared packages & 2 & 460 \\
    \midrule
    \textbf{Total} & \textbf{18} & \textbf{5121} \\
    \bottomrule
  \end{tabular}
\end{table}

\Cref{fig:impl-size} places this total against the reference design. The controller is
roughly one order of magnitude smaller than the reference RTL; the per-subsystem attribution
of that gap, and the scope caveat that must accompany it, are given in
\cref{sec:ref-comparison}.

\begin{figure}[!htb]\centering
\begin{tikzpicture}[font=\footnotesize]
  % bar widths (cm) are pre-scaled from the line counts (\,$\times 0.00018$) to keep the
  % raw counts out of pgfmath, whose intermediate maximum is about 16384.
  \foreach \y/\w/\val/\lab in {
      2/0.92/5121/{This design (RTL)},
      1/6.87/38187/{Reference (RTL only)},
      0/9.03/50160/{Reference (full \texttt{src})}} {
    \fill[blue!55!black] (0,\y) rectangle (\w,{\y+0.55});
    \node[anchor=west] at (\w+0.1,{\y+0.275}) {\val\ lines};
    \node[anchor=east] at (-0.15,{\y+0.275}) {\lab};
  }
  \draw[-{Stealth[length=5pt]}] (0,-0.3) -- (10.2,-0.3) node[anchor=north east]{lines of SystemVerilog};
\end{tikzpicture}
\caption[Implementation-size comparison]{Implementation size of this design against the
  CHIPS Alliance reference. The reference ``RTL only'' bar excludes its generated UVM and
  coverage helper files; the ``full \texttt{src}'' bar includes them.}
\label{fig:impl-size}
\end{figure}

\section{Regression results}
\label{sec:regression}

The full regression --- all \RunnableVseqCount{} runnable scenario sequences --- was run at
the baseline of \cref{tab:results-baseline} with a fixed seed. Every sequence passed with
zero \texttt{UVM\_ERROR} and zero \texttt{UVM\_FATAL} reports.
\Cref{tab:regression-matrix} summarises the outcome by category; the per-sequence log
excerpts are collected in \cref{app:logs}, and the full sequence inventory in
\cref{app:vseq-inventory}.

\begin{table}[!htb]\centering
  \caption{Regression pass/fail matrix by category (seed~1, commit \texttt{047dfdc}).}
  \label{tab:regression-matrix}
  \small
  \begin{tabular}{@{}>{\raggedright\arraybackslash}p{0.52\textwidth} r r@{}}
    \toprule
    \textbf{Category} & \textbf{Tests} & \textbf{Pass} \\
    \midrule
    CSR and DAT access & 14 & 14 \\
    FIFO / queue integrity & 4 & 4 \\
    Bus, PHY, and SCL timing & 11 & 11 \\
    SDR private write & 5 & 5 \\
    SDR private read & 6 & 6 \\
    Immediate data transfer & 4 & 4 \\
    \iic{}-FM legacy transfer & 3 & 3 \\
    CCC (ENEC / DISEC / ENTDAA frame) & 5 & 5 \\
    Dynamic address assignment (DAA) & 6 & 6 \\
    ENTDAA multi-target arbitration & 1 & 1 \\
    Error, abort, and reset handling & 21 & 21 \\
    \midrule
    \textbf{Total} & \textbf{80} & \textbf{80} \\
    \bottomrule
  \end{tabular}
\end{table}

\section{Waveform evidence}
\label{sec:waveforms}

Waveform databases were captured for the four canonical transaction classes by re-running
the corresponding sequences with waveform dumping enabled; each database is written to
\texttt{waves/<sequence>/waves.shm}. \Cref{fig:wave-write,fig:wave-read,fig:wave-imm,%
fig:wave-ccc} present the annotated captures, cropped to the first few microseconds and
overlaid with markers on the START, address ACK, T-Bit, open-drain-to-push-pull handover,
and STOP events that define each protocol phase.

\figphl{fig:wave-write}{SDR private write, annotated}
  {START, \texttt{7E}/dynamic-address ACK, OD$\rightarrow$PP handover, data byte + odd-parity
   T-Bit, STOP}{\texttt{waves/i3c\_write\_vseq/waves.shm}, seed~1}{SimVision export}

\figphl{fig:wave-read}{SDR private read, annotated}
  {START, address ACK, data byte + end-of-data T-Bit, target-driven data phase, STOP}
  {\texttt{waves/i3c\_read\_vseq/waves.shm}, seed~1}{SimVision export}

\figphl{fig:wave-imm}{Immediate data transfer, annotated}
  {descriptor-embedded data bytes issued without a separate TX-FIFO fetch}
  {\texttt{waves/i3c\_imm\_vseq/waves.shm}, seed~1}{SimVision export}

\figphl{fig:wave-ccc}{Broadcast CCC (ENEC) and ENTDAA frame, annotated}
  {broadcast \texttt{7E} header, CCC opcode, and the ENTDAA per-target ID + dynamic-address
   loop}{\texttt{waves/i3c\_ccc\_broadcast\_enec\_vseq/waves.shm} and
   \texttt{waves/i3c\_ccc\_entdaa\_opening\_frame\_vseq/waves.shm}, seed~1}{SimVision export}

\section{Functional coverage}
\label{sec:funccov}

Functional coverage was collected over the same \RunnableVseqCount{} sequences, merged into a
single database, and post-processed by the project's own extractor into the committed
\texttt{coverage\_report.txt} (\cref{app:logs}). \Cref{tab:cov-summary} gives the headline
result: 78.8\,\% of coverage bins are hit, and 83 of the 97 defined coverpoints (85.6\,\%)
are fully closed.

\begin{table}[!htb]\centering
  \caption{Functional-coverage summary (merged over \RunnableVseqCount{} sequences,
    \texttt{coverage\_report.txt}).}
  \label{tab:cov-summary}
  \small
  \begin{tabular}{@{}>{\raggedright\arraybackslash}p{0.60\textwidth} r@{}}
    \toprule
    \textbf{Metric} & \textbf{Value} \\
    \midrule
    Sequences analysed & 80 \\
    Coverage bins defined & 1147 \\
    Coverage bins hit & 904 \\
    Overall bin coverage & 78.8\,\% \\
    Coverpoints defined & 97 \\
    Coverpoints fully closed (100\,\%) & 83 (85.6\,\%) \\
    Cross covergroups defined & 36 \\
    Cross covergroups fully closed & 19 \\
    \bottomrule
  \end{tabular}
\end{table}

The gap between 85.6\,\% closed coverpoints and 78.8\,\% overall bins is informative: the
individual protocol coverpoints are almost entirely closed, and the residual comes from two
identifiable sources rather than from untested behaviour --- chiefly the automatic cross
covergroups, of which only 19 of 36 close. \Cref{tab:cov-gaps} lists the representative
unclosed coverpoints with the reason each is open.

\begin{table}[!htb]\centering
  \caption{Representative unclosed coverpoints and the reason each remains open.}
  \label{tab:cov-gaps}
  \small
  \begin{tabular}{@{}>{\raggedright\arraybackslash}p{0.34\textwidth} c
                    >{\raggedright\arraybackslash}p{0.42\textwidth}@{}}
    \toprule
    \textbf{Coverpoint} & \textbf{Hit/Total} & \textbf{Reason open} \\
    \midrule
    Automatic cross groups & 24--96\,\% & structurally unreachable field combinations in a
      single-controller directed environment \\
    \texttt{cp\_cmd\_dat\_idx} & 9/32 & directed tests exercise a subset of the 32 DAT slots \\
    \texttt{cp\_cmd\_dev\_count} & 4/16 & ENTDAA device-count breadth beyond directed scenarios \\
    \texttt{cp\_identity\_pattern} & 1/4 & multi-target arbitration PID-pattern breadth \\
    \texttt{cp\_round\_count}, \texttt{cp\_target\_count} & 1/3 & multi-round, multi-target
      DAA breadth \\
    \texttt{cp\_first\_loss\_phase} & 2/3 & arbitration loss-phase breadth \\
    \texttt{cp\_cmd\_dbp}, \texttt{cp\_ccc\_event\_t\_bit}, \texttt{cp\_ccc\_header\_nack}
      & 1/2 & individually rare stimulus conditions \\
    \bottomrule
  \end{tabular}
\end{table}

The first source is the automatically generated cross covergroups, whose Cartesian
combinations include field pairings that cannot co-occur for a single Active Controller. The
second is a small set of \emph{breadth} coverpoints --- all 32 DAT indices, all 16 ENTDAA
device counts, and the multi-target arbitration permutations --- that lie at or beyond the
directed scope of a single-author controller. No protocol \emph{mode} is unexercised; the
open bins are combinatorial breadth, not missing functionality. Closing them with
constrained-random sweeping is identified as future work in \cref{ch:conclusion}.

\section{Comparison with CHIPS Alliance reference}
\label{sec:ref-comparison}

\Cref{tab:ref-comparison} attributes the size difference of \cref{fig:impl-size} to
individual subsystems. On synthesisable RTL, the controller is 5121 lines against the
reference's 38187 --- an 87\,\% reduction; against the full reference \texttt{src} tree
(50160 lines, including its generated UVM/coverage helpers) the reduction is 90\,\%.

\begin{table}[!htb]\centering
  \caption{Per-subsystem line count, this design versus the reference (RTL only).}
  \label{tab:ref-comparison}
  \small
  \begin{tabular}{@{}>{\raggedright\arraybackslash}p{0.40\textwidth} r r
                    >{\raggedright\arraybackslash}p{0.24\textwidth}@{}}
    \toprule
    \textbf{Subsystem} & \textbf{This} & \textbf{Ref.} & \textbf{Note} \\
    \midrule
    Control FSMs / datapath & 3997 & 13842 & ref.\ adds IBI, HDR, secondary controller,
      target mode \\
    CSR + DAT & 368 & 14296 & ref.\ CSR is PeakRDL-generated; this design is hand-written \\
    HCI queues & 241 & 2665 & ref.\ full HCI/PIO/DMA \\
    PHY & 55 & 119 & \\
    Top level + packages & 460 & 2304 & ref.\ wrapper, interrupt, defines \\
    OCP recovery & --- & 3057 & out of this design's scope \\
    AXI / memory libraries & --- & 1904 & out of this design's scope \\
    \midrule
    \textbf{Total (RTL)} & \textbf{5121} & \textbf{38187} & 87\,\% reduction \\
    \bottomrule
  \end{tabular}
\end{table}

This reduction must be read with its scope caveat. It is not a claim that the same design was
written more densely. Most of the difference is structural: the reference implements IBI,
HDR, multi-master/secondary-controller and target modes, an OCP recovery block, and a full
AXI-based host interface, none of which are in this design's scope
(\cref{ch:background-requirements}); and its CSR block is machine-generated from a register description,
whereas this design hand-writes a 26-register file with a 32-entry DAT. What the comparison
does establish is the thesis contribution stated in \cref{ch:intro-full}: a study-scale Active
Controller that still implements all \FlowStateCount{} \texttt{flow\_active} states and
retains SDR, \iic{}-FM, and ENTDAA, in roughly a tenth of the reference's RTL. The verified
feature evidence for that claim is the green regression of \cref{tab:regression-matrix} and
the coverage of \cref{tab:cov-summary}.

\section{Evaluation against requirements}
\label{sec:req-eval}

\Cref{tab:req-eval} closes the loop against the requirements of \cref{ch:background-requirements}, marking
each as met, partially met, or not demonstrated, with the evidence from this chapter.

\begin{table}[!htb]\centering
  \caption{Evaluation against the requirements of \cref{ch:background-requirements}.}
  \label{tab:req-eval}
  \small
  \begin{tabular}{@{}>{\raggedright\arraybackslash}p{0.40\textwidth}
                    >{\raggedright\arraybackslash}p{0.16\textwidth}
                    >{\raggedright\arraybackslash}p{0.36\textwidth}@{}}
    \toprule
    \textbf{Requirement} & \textbf{Status} & \textbf{Evidence} \\
    \midrule
    SDR private write / read & Met & regression 11/11 (SDR) + scoreboard data/T-Bit checks \\
    Immediate data transfer & Met & regression 4/4 (immediate) \\
    \iic{}-FM legacy write / read & Met & regression 3/3 (\iic{}) \\
    ENTDAA dynamic addressing & Met & regression 6/6 (DAA) + 1/1 arbitration \\
    ENEC / DISEC CCC & Met & regression 5/5 (CCC) \\
    Error, abort, and reset handling & Met & regression 21/21 + \texttt{ERR\_STATUS} checks \\
    CSR + 32-entry DAT & Met & regression 14/14 (CSR) + register model \\
    FIFO / queue integrity & Met & regression 4/4 + \texttt{sync\_fifo} assertions \\
    Bus, PHY, and SCL timing & Met & regression 11/11 + bound timing assertions \\
    SDR \SI{12.5}{\mega\hertz} / \iic{} \SI{400}{\kilo\hertz} & Met (functional) &
      bus timing sequences + SCL assertions \\
    System clock $\ge \SI{333}{\mega\hertz}$ & Not demonstrated & no synthesis run;
      deferred to future work (\cref{ch:conclusion}) \\
    Functional-coverage closure & Partially met & 78.6\,\% bins, 80.5\,\% coverpoints;
      remainder is out-of-scope breadth (\cref{tab:cov-gaps}) \\
    \bottomrule
  \end{tabular}
\end{table}

Every functional and interface requirement is met with regression and, where applicable,
assertion evidence. The two items short of ``met'' are bounded and understood: the
$\ge \SI{333}{\mega\hertz}$ timing target is a synthesis property that the simulation-only
flow of this thesis does not measure, and coverage closure is partial only in the
combinatorial-breadth sense established in \cref{sec:funccov}. Both are carried forward as
future work in \cref{ch:conclusion}.
